\renewcommand{\figurename}{Handout}

\newcounter{handoutCount}
\setcounter{handoutCount}{0}

\makeatletter
\NewDocumentCommand\defineTextHandout{m m m}{
    \defineHandout{#1}{#2}{\begin{tcolorbox}#3\end{tcolorbox}}
}
\NewDocumentCommand\defineImageHandout{m m m}{
    \defineHandout{#1}{#2}{\image{#3}}
}
\NewDocumentCommand\defineHandout{m m m}{

    \global\expandafter\gdef\csname handout@#1\endcsname{
        #3
        \noindent\begin{minipage}{\linewidth}
            \captionsetup{labelformat=empty}
            \captionof{figure}{Handout \handoutNum{#1}: #2}
        \end{minipage}
    }
}
\NewDocumentCommand\handout{m}{%

    \ifcsname handout@#1\endcsname%
        \csname handout@#1\endcsname%
    \else%
        \PackageError{handout}{Handout `#1' not defined}{}%
    \fi%
}

\NewDocumentCommand\handoutNum{m}{%
    \ifcsname handoutNum@#1\endcsname%
    \else%
        \stepcounter{handoutCount}%
        \global\expandafter\edef\csname handoutNum@#1\endcsname {\thehandoutCount}
    \fi%
    \ifcsname handoutNum@#1\endcsname%
        \csname handoutNum@#1\endcsname%
    \else%
        \PackageError{handout}{Handout number for `#1' not defined}{}%
    \fi%
}

\NewDocumentCommand\refhandout{m}{%
    (\textbf{Handout \handoutNum{#1}})
}

% \loadHandouts[<folder>]
% Scans <folder> (default: ./handouts) for .tex files and \input each one.
% Also accepts a custom \input prefix when the folder path contains one level
% of nesting that should be used as-is (e.g. \loadHandouts[./new_handouts/handouts]).
\NewDocumentCommand\loadHandouts{O{./handouts}}{%
    \directlua{require("parser").emit_handout_inputs([[#1]])}%
}
\makeatother
